\section{Context 数据模型：AI 社交到底要存什么}

要让分身像用户，仅有资料卡远远不够。AI 社交需要的是多层 context：稳定身份、长期偏好、近期状态、表达风格、社交目标、关系图谱、隐私边界和反馈历史。不同层的时效性、敏感度和授权方式都不同，不能混成一个“记忆”字段。

\subsection{Context 分层}

\noindent\begin{longtable}{p{0.22\textwidth}p{0.34\textwidth}p{0.35\textwidth}}
\toprule
层级 & 内容 & 产品处理 \\
\midrule
身份层 & 基本身份、职业、城市、公开 profile & 可公开、可编辑、低敏感。 \\
偏好层 & 兴趣、审美、喜欢/讨厌的表达 & 用于推荐和语气生成。 \\
关系层 & 朋友、关注、历史互动、弱关系 & 高价值也高敏感，需要权限。 \\
状态层 & 近期情绪、目标、正在做的事 & 时效性强，需要过期机制。 \\
表达层 & 写作风格、幽默、观点、边界 & 决定分身像不像用户。 \\
反馈层 & 用户认可/否定分身行为 & 训练分身和调整策略的核心。 \\
\bottomrule
\end{longtable}

\begin{importantbox}{Context 不是越多越好，而是越分层越好}
不同 context 层级的敏感度、时效性和用途不同。AI 社交产品如果不分层管理，就会同时遇到隐私风险和体验混乱。
\end{importantbox}

\subsection{Context 的生命周期}

Context 应有生命周期：采集、解释、使用、反馈、过期、删除。比如“我今天想认识做 AI 产品的人”是短期目标，不应永久固化；“我讨厌冒犯式幽默”是长期偏好，可以长期保留；“某段私聊内容”则可能只能本地使用，不应进入公共匹配。

\begin{warningbox}{记忆污染问题}
如果分身把一次临时表达当成长期偏好，或者把调侃当成价值观，就会产生记忆污染。AI 社交产品必须允许用户纠正和删除 context，否则越用越不像自己。
\end{warningbox}

\subsection{本章小结}

Context 是 AI 社交的底层数据模型。它需要分层、授权、过期和反馈，而不是简单堆积聊天记录。

